\chapter{Library Dependencies}
\label{ch:libdeps}

\newthought{Some workflows need shared libraries} built and installed before any calculator or
opera can run---think HepMC, Pythia8, or Delphes. The optional \yamlkey{LibDeps} block declares
those libraries inside the task card, so the setup is reproducible and part of the project rather
than a hand-run script.

\section{What LibDeps is for}
\label{sec:lib-why}

\yamlkey{LibDeps} answers three questions about supporting libraries:

\begin{enumerate}
  \item \textbf{Where} should they live? --- a shared root directory.
  \item \textbf{How} is each one installed? --- a reproducible command recipe.
  \item \textbf{In what order}? --- explicit prerequisites.
\end{enumerate}

\noindent These libraries are usually neither calculators nor operas; they provide build-time or
run-time support for the tools that are.

\section{The two layers}
\label{sec:lib-layers}

\yamlkey{LibDeps} has a shared top level and a list of \yamlkey{Modules}, one per library:

\begin{yamlwin}
LibDeps:
  path: "&J/deps/library"
  make_paraller: 16
  Modules:
    - name: "HepMC"
      required_modules: []
      installed: False
      installation:
        path: "&J/deps/library/HepMC"
        source: "&J/deps/library/source/HepMC-2.06.09.tar.gz"
        commands:
          - "..."
\end{yamlwin}

\begin{description}[style=nextline, leftmargin=1.2em]
  \item[\yamlkey{path}] the shared root for all dependency files; reusable in commands as
        \code{\$\{LibDeps:path\}}.
  \item[\yamlkey{make\_paraller}] default build parallelism; reusable as
        \code{make -j\$\{LibDeps:make\_paraller\}}.
\end{description}

\noindent Each module has a \yamlkey{name}, a \yamlkey{required\_modules} list, an
\yamlkey{installed} flag, and an \yamlkey{installation} section (\yamlkey{path},
\yamlkey{source}, and ordered \yamlkey{commands}).

\section{Variable substitution in commands}
\label{sec:lib-vars}

Inside \yamlkey{installation.commands} you can use these placeholders, which keep the recipe
readable and portable:

\begin{description}[style=nextline, leftmargin=1.2em]
  \item[\code{\$\{LibDeps:path\}}] the shared root;
  \item[\code{\$\{LibDeps:make\_paraller\}}] the shared build parallelism;
  \item[\code{\$\{path\}}] this module's \yamlkey{installation.path};
  \item[\code{\$\{source\}}] this module's \yamlkey{installation.source}.
\end{description}

\section{Installation order}
\label{sec:lib-order}

Declare prerequisites with \yamlkey{required\_modules}---for example, Pythia8 needs HepMC:

\begin{yamlwin}
- name: "Pythia8"
  required_modules: ["HepMC"]
\end{yamlwin}

\begin{jwarn}
In the current version, \yamlkey{required\_modules} on a \yamlkey{LibDeps} module is recorded as
\emph{metadata}---it documents intent but is not auto-enforced. Install libraries in a sensible
order yourself: a dependency must be built before the library that configures against it (HepMC
before Pythia8).
\end{jwarn}

\section{A worked example: HepMC then Pythia8}
\label{sec:lib-example}

\begin{yamlwin}[LibDeps block]
LibDeps:
  path: "&J/deps/library"
  make_paraller: 16
  Modules:
    - name: "HepMC"
      required_modules: []
      installed: False
      installation:
        path: "&J/deps/library/HepMC"
        source: "&J/deps/library/source/HepMC-2.06.09.tar.gz"
        commands:
          - "cd ${LibDeps:path}"
          - "rm -rf HepMC*"
          - "cp ${source} ./"
          - "tar -xzf HepMC-2.06.09.tar.gz"
          - "cd ${LibDeps:path}/HepMC-2.06.09"
          - "./configure --with-momentum=GEV --with-length=MM --prefix=${path}"
          - "make -j${LibDeps:make_paraller}"
          - "make install"
    - name: "Pythia8"
      required_modules: ["HepMC"]
      installed: False
      installation:
        path: "&J/deps/library/Pythia8"
        source: "&J/deps/library/source/pythia8230.tgz"
        commands:
          - "cd ${LibDeps:path}/Source"
          - "tar -zxvf ${source}"
          - "cd ${LibDeps:path}/Source/pythia8230"
          - "./configure --with-hepmc2=${LibDeps:path}/HepMC --prefix=${path}"
          - "make -j${LibDeps:make_paraller}"
          - "make install"
\end{yamlwin}

\noindent Pythia8's \code{-{}-with-hepmc2} flag points at HepMC's installed
\yamlkey{path}---which is exactly why HepMC must be built first.

\section{What Jarvis does at runtime}
\label{sec:lib-runtime}

When \Jarvis\ loads a card with \yamlkey{LibDeps}, it resolves the paths and placeholders, then
runs each module's \yamlkey{commands} in order, streaming their output to a per-module install
log. It also writes a small installation record per module.

\begin{jnote}
If \Jarvis\ finds a matching installation record, it treats the library as already installed and
\emph{asks} whether to reinstall---it does not silently rebuild. You can skip library
installation entirely for a run with \cmd{-{}-skip-library-installation}
(Chapter~\ref{ch:cli}).
\end{jnote}

\section{Summary}
\label{sec:lib-summary}

\yamlkey{LibDeps} turns one-time library setup into a reproducible, project-local recipe: a shared
\yamlkey{path}, one \yamlkey{Modules} entry per library with ordered \yamlkey{commands}, and
\code{\$\{...\}} placeholders for portability. Build prerequisites first, and let the card
document your environment.
